% Options for packages loaded elsewhere
\PassOptionsToPackage{unicode}{hyperref}
\PassOptionsToPackage{hyphens}{url}
\documentclass[
]{book}
\usepackage{xcolor}
\usepackage{amsmath,amssymb}
\setcounter{secnumdepth}{5}
\usepackage{iftex}
\ifPDFTeX
  \usepackage[T1]{fontenc}
  \usepackage[utf8]{inputenc}
  \usepackage{textcomp} % provide euro and other symbols
\else % if luatex or xetex
  \usepackage{unicode-math} % this also loads fontspec
  \defaultfontfeatures{Scale=MatchLowercase}
  \defaultfontfeatures[\rmfamily]{Ligatures=TeX,Scale=1}
\fi
\usepackage{lmodern}
\ifPDFTeX\else
  % xetex/luatex font selection
\fi
% Use upquote if available, for straight quotes in verbatim environments
\IfFileExists{upquote.sty}{\usepackage{upquote}}{}
\IfFileExists{microtype.sty}{% use microtype if available
  \usepackage[]{microtype}
  \UseMicrotypeSet[protrusion]{basicmath} % disable protrusion for tt fonts
}{}
\makeatletter
\@ifundefined{KOMAClassName}{% if non-KOMA class
  \IfFileExists{parskip.sty}{%
    \usepackage{parskip}
  }{% else
    \setlength{\parindent}{0pt}
    \setlength{\parskip}{6pt plus 2pt minus 1pt}}
}{% if KOMA class
  \KOMAoptions{parskip=half}}
\makeatother
\usepackage{longtable,booktabs,array}
\newcounter{none} % for unnumbered tables
\usepackage{calc} % for calculating minipage widths
% Correct order of tables after \paragraph or \subparagraph
\usepackage{etoolbox}
\makeatletter
\patchcmd\longtable{\par}{\if@noskipsec\mbox{}\fi\par}{}{}
\makeatother
% Allow footnotes in longtable head/foot
\IfFileExists{footnotehyper.sty}{\usepackage{footnotehyper}}{\usepackage{footnote}}
\makesavenoteenv{longtable}
\usepackage{graphicx}
\makeatletter
\newsavebox\pandoc@box
\newcommand*\pandocbounded[1]{% scales image to fit in text height/width
  \sbox\pandoc@box{#1}%
  \Gscale@div\@tempa{\textheight}{\dimexpr\ht\pandoc@box+\dp\pandoc@box\relax}%
  \Gscale@div\@tempb{\linewidth}{\wd\pandoc@box}%
  \ifdim\@tempb\p@<\@tempa\p@\let\@tempa\@tempb\fi% select the smaller of both
  \ifdim\@tempa\p@<\p@\scalebox{\@tempa}{\usebox\pandoc@box}%
  \else\usebox{\pandoc@box}%
  \fi%
}
% Set default figure placement to htbp
\def\fps@figure{htbp}
\makeatother
\setlength{\emergencystretch}{3em} % prevent overfull lines
\providecommand{\tightlist}{%
  \setlength{\itemsep}{0pt}\setlength{\parskip}{0pt}}
\usepackage[]{natbib}
\bibliographystyle{plainnat}
\usepackage{booktabs}

\usepackage{color}
\usepackage{framed}
\setlength{\fboxsep}{.8em}

% These colours were manually entered, they shouldn't matter unless you want pdf output

\newenvironment{redbox}{
  \definecolor{shadecolor}{RGB}{243, 154, 157}
  \color{white}
  \begin{shaded}}
 {\end{shaded}}

\newenvironment{bluebox}{
  \definecolor{shadecolor}{RGB}{172, 210, 237}
  \color{white}
  \begin{shaded}}
 {\end{shaded}}

\newenvironment{greenbox}{
  \definecolor{shadecolor}{RGB}{141, 181, 128}
  \color{white}
  \begin{shaded}}
 {\end{shaded}}
\usepackage{bookmark}
\IfFileExists{xurl.sty}{\usepackage{xurl}}{} % add URL line breaks if available
\urlstyle{same}
\hypersetup{
  pdftitle={DRAC Primer Session 2026},
  hidelinks,
  pdfcreator={LaTeX via pandoc}}

\title{DRAC Primer Session 2026}
\author{}
\date{\vspace{-2.5em}}

\begin{document}
\maketitle

{
\setcounter{tocdepth}{1}
\tableofcontents
}
\chapter{DRAC Primer Info}\label{drac-primer-info}

Coming soon!

\subsection{Data download}\label{data-download}

Coming soon!

\subsection{Schedule}\label{schedule}

\textbf{High-Throughput Genomics DRAC Pre-Session 2026}: Wednesday, October 28, 2026

{\def\LTcaptype{none} % do not increment counter
\begin{longtable}[]{@{}cc@{}}
\toprule\noalign{}
(PDT) & Module \\
\midrule\noalign{}
\endhead
\bottomrule\noalign{}
\endlastfoot
11:00 & Welcome Introduction \\
11:10 & Module 1: Logging into DRAC \\
12:10 & Short Break \\
12:15 & Module 2: Introduction to the Command Line \\
12:35 & Module 3: File Manipulation \\
13:05 & Long Break \\
14:05 & Module 4: Searching and Sorting Through Files \\
15:00 & Short Break \\
15:05 & Module 5: Shell scripts \\
16:00 & Finished \\
\end{longtable}
}

\chapter{Meet Your Faculty}\label{meet-your-faculty}

\subsection{Instructors}\label{instructors}

\subsubsection{Jiarui Li}\label{jiarui-li}

\begin{quote}
JOB TITLE
INSTITUTION
LOCATION

--- CONTACT INFO, IF PROVIDED
\end{quote}

BIO GOES HERE--\textgreater{}

\subsection{Facilitator}\label{facilitator}

\subsubsection{Tian Rabbani}\label{tian-rabbani}

\begin{quote}
Regional Coordinator, CBH British Columbia
Canadian Bioinformatics Hub, Simon Fraser University
Vancouver, BC, Canada

--- \href{mailto:bc@bioinformatics.ca}{\nolinkurl{bc@bioinformatics.ca}}
\end{quote}

Tian completed her Master of Public Health in 2022 from Simon Fraser University and is currently a researcher and project lead at the Center for Infectious Disease Genomics and One Health (CIDGOH). She has been extensively involved in researching SARS-CoV-2 data-sharing through the Canadian COVID-19 Genomics Network (CanCOGeN), alongside leading efforts to understand metadata reporting in genomic epidemiology studies for outbreak and surveillance investigation. Her background in public health and knowledge in contextual data reporting intersect harmoniously to help inform current and future emergency response initiatives.

As British Columbia's regional coordinator for the Canadian Bioinformatics Hub, Tian is passionate about and dedicated to connecting bioinformatics communities and groups across the province. She aims to grow, develop, and expand data science and bioinformatics tools and training materials for prospective or new professionals, as well as for those looking to advance in their current careers. Tian believes accessibility, access, and availability of these tools and training opportunities are the cornerstone of growing our bioinformatics communities across the region.

\bibliography{book.bib,packages.bib}

\end{document}
